\chapter{GENERAL FRAMEWORK}
\begin{spacing}{1.2}
\minitoc
\thispagestyle{MyStyle}
\end{spacing}
\newpage

\section{INTRODUCTION}
This chapter presents the general framework of our final-year project. We begin by introducing the host organization CLINISYS, followed by a critical analysis of the existing system. We then define the core problem statement, present the proposed SIPMS solution, and describe the Scrum agile methodology adopted for the development process.

\section{PRESENTATION OF THE HOST ORGANIZATION}

\subsection{About CLINISYS}
CLINISYS is a technology company specializing in the development of innovative digital solutions for industrial and academic sectors. With deep expertise in software engineering, the company supports its clients in their digital transformation journeys by delivering custom-built systems that combine performance, security, and usability.

\par
As part of its commitment to talent development and innovation, CLINISYS welcomes a growing number of interns from higher education institutions each year. This ongoing engagement with academic talent has highlighted the urgent need for a structured digital tool capable of managing the full internship lifecycle efficiently.

\subsection{Key Activity Areas}
\renewcommand{\labelitemi}{\tiny$\bullet$}
\begin{itemize}[leftmargin=2cm, topsep=0pt]
    \begin{spacing}{1.25}
        \item Custom web and mobile application development
        \item Digital transformation consulting and process optimization
        \item Artificial intelligence and data analytics integration
        \item Enterprise information system management and hosting
    \end{spacing}
\end{itemize}

\section{ANALYSIS OF THE EXISTING SYSTEM}

\subsection{Description of the Current System}
Prior to the development of SIPMS, CLINISYS managed its internship and project workflows using entirely manual processes. Candidate registration was performed on paper forms, communication relied on unstructured emails, and progress tracking depended on Excel spreadsheets maintained manually by reception staff. Document management — including CVs, cover letters, and academic transcripts — was disorganized and difficult to query.

\subsection{Identified Limitations}
The analysis of the existing system revealed several critical dysfunctions, summarized in the table below:

\begin{table}[H]
\centering
\begin{tabular}{|p{4.5cm}|p{8.5cm}|}
\hline
\rowcolor{gray!20}
\textbf{Identified Problem} & \textbf{Organizational Impact} \\
\hline
Manual candidate registration & High risk of data entry errors, loss of records, slow onboarding process \\
\hline
No competency evaluation system & Inability to objectively assess candidates' technical skills \\
\hline
Non-automated communication & Significant delays in sending login credentials and notifications \\
\hline
No year-based filtering & Difficulty retrieving and managing candidates from a given intake year \\
\hline
Absent document management & CVs and supporting documents stored in an unstructured manner \\
\hline
Manual project assignment & Subjective, time-consuming process with no defined relevance criteria \\
\hline
\end{tabular}
\caption{Summary of existing system limitations}
\end{table}

\section{PROBLEM STATEMENT}
Given the shortcomings of the current system, CLINISYS expressed the need for a centralized digital solution capable of handling the complete internship candidate lifecycle. The core problem statement is formulated as follows:

\par
\textit{``How can a smart web platform be designed and developed to automate and centralize the management of internship candidates — from registration to supervised project assignment — while ensuring data security, action traceability, and an optimal user experience for all stakeholders involved?''}

\section{PROPOSED SOLUTION}

\subsection{The SIPMS Platform}
In response to this problem, we propose the development of the \textbf{Smart Internship and Project Management System (SIPMS)}, a full-stack, multi-role web platform. The solution provides:

\begin{itemize}[leftmargin=2cm, topsep=0pt]
    \begin{spacing}{1.25}
        \item A \textbf{Reception Panel} enabling receptionists to register candidates, upload their documents (CV, cover letter, academic transcript, ID card), and filter candidates by internship year.
        \item An \textbf{automated quiz system} for objective, role-specific evaluation of candidates' technical competencies.
        \item An \textbf{AI engine} for project ranking, CV analysis, roadmap generation, and intelligent matching between candidates and available supervisors.
        \item An \textbf{automated email system} that instantly sends login credentials upon candidate account creation.
        \item \textbf{Role-differentiated dashboards} tailored to each user type: Administrator, Manager, Receptionist, Candidate, and Supervisor.
        \item \textbf{Complete application lifecycle management}, tracking each candidature from submission through review, quiz evaluation, AI matching, interview, and final project assignment.
        \item A \textbf{CV upload and AI analysis module} allowing candidates to upload their CV and receive automated skill extraction, project matching recommendations, and a personalized 4-week roadmap.
    \end{spacing}
\end{itemize}

\subsection{Technology Stack Overview}
\begin{table}[H]
\centering
\begin{tabular}{|p{3.5cm}|p{4cm}|p{5.5cm}|}
\hline
\rowcolor{gray!20}
\textbf{Layer} & \textbf{Technology} & \textbf{Role} \\
\hline
Frontend & React (Vite) + CSS & Reactive single-page application \\
\hline
Backend & Spring Boot (Java 17) & RESTful API, business logic \\
\hline
Security & Spring Security + JWT & Stateless authentication, RBAC \\
\hline
Database & MySQL 8 & Relational data persistence \\
\hline
Email & JavaMail (Gmail SMTP) & Automated transactional emails \\
\hline
File Storage & Spring Multipart & CV and document upload \\
\hline
AI Module & Custom scoring engine & Project ranking, candidate matching \\
\hline
\end{tabular}
\caption{SIPMS technology stack}
\end{table}

\section{DEVELOPMENT METHODOLOGY}

\subsection{Choice of Methodology: Scrum}
We adopted the \textbf{Scrum} agile framework for this project. This choice is justified by the iterative nature of the requirements, the need to deliver functional increments regularly, and the flexibility required to accommodate changes during development.

\subsection{Scrum Framework}
Scrum organizes development into short iterations called \textbf{Sprints}, each lasting two to three weeks. Every sprint produces a potentially shippable product increment. The key Scrum artifacts used in this project are:

\begin{itemize}[leftmargin=2cm, topsep=0pt]
    \begin{spacing}{1.25}
        \item \textbf{Product Backlog}: A prioritized list of all features to be developed.
        \item \textbf{Sprint Backlog}: The subset of the Product Backlog selected for a given sprint.
        \item \textbf{Increment}: A functional version of the software delivered at the end of each sprint.
    \end{spacing}
\end{itemize}

\subsection{Product Backlog}
The Product Backlog is the single source of truth for all system requirements. It contains 34 user stories prioritized by their business value.

\begin{longtable}{|c|p{2.5cm}|p{8cm}|c|}
\hline
\rowcolor{gray!20}
\textbf{ID} & \textbf{Module} & \textbf{User Story} & \textbf{Priority} \\
\hline
\endfirsthead

\hline
\rowcolor{gray!20}
\textbf{ID} & \textbf{Module} & \textbf{User Story} & \textbf{Priority} \\
\hline
\endhead

US-01 & Auth & As a User, I want to login with role-based access (RBAC) to access only features for my profile. & High \\
\hline
US-02 & Auth & As an Admin/Receptionist, I want to create user accounts with auto-generated or temporary passwords. & High \\
\hline
US-03 & Auth & As a User, I must change my temporary password on first login for security. & High \\
\hline
US-04 & Management & As a Receptionist, I want to register physical dossiers in the system. & High \\
\hline
US-05 & Management & As an Admin/Manager, I want to manage users (create, modify, delete, toggle active). & High \\
\hline
US-06 & Management & As an Admin/Manager, I want to manage supervisors to assign them to projects. & High \\
\hline
US-07 & Candidate & As a Candidate, I want to submit an online application for a project to be evaluated by the system. & High \\
\hline
US-08 & AI System & As a Manager, I want the AI to analyze and rank submitted projects to validate the most relevant ones. & High \\
\hline
US-09 & AI System & As a Manager, I want the AI to suggest the most suitable supervisor. & Medium \\
\hline
US-10 & Evaluation & As a Candidate, I want to take a timed technical quiz to prove my skills. & High \\
\hline
US-11 & Evaluation & As a System, I want to auto-correct quizzes and generate scores. & High \\
\hline
US-12 & Evaluation & As a Candidate, I want to take a quiz specific to my specialty within 48 hours of assignment. & High \\
\hline
US-13 & Management & As a Manager, I want to set candidate application status with valid state transitions. & High \\
\hline
US-14 & Notification & As a User, I want to receive automatic notifications at each key step. & Medium \\
\hline
US-15 & Dashboard & As a User, I want to be redirected to my specific panel based on my role after login. & High \\
\hline
US-16 & Internship File & As a Receptionist, I want to add per-year internship files with optional document upload. & High \\
\hline
US-17 & Internship File & As a Receptionist/Manager, I want to view all internship files across candidates in a dedicated tab. & High \\
\hline
US-18 & Quiz & As a Manager, I want to approve a candidate and send a quiz. & High \\
\hline
US-19 & Interview & As a Manager, I want to schedule interviews for candidates who passed the quiz. & High \\
\hline
US-20 & Interview & As a Manager, I want to record interview results and update the interview status. & High \\
\hline
US-21 & Interview & As a Manager, I want to view all interviews in a dedicated interview management page. & High \\
\hline
US-22 & AI System & As a Candidate, I want the AI to analyze my CV and generate a 4-week project roadmap. & Medium \\
\hline
US-23 & AI System & As a Candidate/Manager, I want the AI to match my CV against available projects and rank the top 3 best fits. & High \\
\hline
US-24 & Management & As an Admin/Manager, I want to manage projects to maintain the project catalog. & High \\
\hline
US-25 & Management & As a Candidate, I want to propose a project idea to be evaluated by the AI and reviewed by the manager. & Medium \\
\hline
US-26 & Application & As a Receptionist, I want to register a physical application with candidate linked to a specific project. & High \\
\hline
US-27 & Application & As a Manager, I want to assign a supervisor to an application considering supervisor capacity constraints. & High \\
\hline
US-28 & Dashboard & As a Manager, I want to view dashboard statistics for monitoring. & Medium \\
\hline
US-29 & Notifications & As a User, I want to view, mark as read, and dismiss in-app notifications. & Low \\
\hline
US-30 & Quiz & As an Admin/Manager, I want to create quizzes with multiple-choice questions. & High \\
\hline
US-31 & Reception & As a Receptionist, I want a dedicated reception panel to manage candidate creation. & High \\
\hline
US-32 & Candidate & As a Candidate, I want to view and edit my profile and track my application status. & Medium \\
\hline
US-33 & Audit & As an Admin, I want to view audit logs tracking all system actions for compliance. & Low \\
\hline
US-34 & AI System & As a Candidate, I want to upload my CV document and have the AI extract skills and match me to suitable projects. & High \\
\hline
\caption{Product Backlog for the SIPMS platform}
\label{tab:product-backlog}
\end{longtable}

\subsection{Sprint Goals}
Each sprint was driven by a clearly defined goal that aligned with the overall project roadmap:

\renewcommand{\labelitemi}{\tiny$\bullet$}
\begin{itemize}[leftmargin=2cm, topsep=4pt]
    \begin{spacing}{1.25}
        \item \textbf{Sprint 1 — Authentication \& User Management}: Establish a secure access layer with role-based dashboards. Goal: enable login/logout with JWT, enforce RBAC, and provide personalized dashboards for Admin, Manager, Receptionist, and Candidate roles.
        \item \textbf{Sprint 2 — Reception Module}: Digitize the physical candidate registration process. Goal: allow receptionists to register candidates, upload documents, filter by year, and send automated welcome emails.
        \item \textbf{Sprint 3 — Quiz \& Project Module}: Implement objective technical evaluation and project configuration. Goal: enable managers to create quizzes, configure projects, and the system to auto-correct quizzes and trigger real-time notifications to managers.
        \item \textbf{Sprint 4 — Interview, AI \& Assignment Module}: Enhance decision-making with artificial intelligence and structured interviews. Goal: enforce quiz-pass gateways for interviews, rank submitted projects using AI, match candidates to supervisors, and finalize the assignment workflow.
    \end{spacing}
\end{itemize}

\subsection{Validation per Sprint}
At the end of each sprint, the delivered increment was validated through a structured process to ensure quality and alignment with requirements:

\renewcommand{\labelitemi}{\tiny$\bullet$}
\begin{itemize}[leftmargin=2cm, topsep=4pt]
    \begin{spacing}{1.25}
        \item \textbf{Sprint Review (Demo)}: A live demonstration of the working software was presented to the project supervisor and CLINISYS stakeholders. Each user story completed during the sprint was executed in real time to confirm functionality.
        \item \textbf{Acceptance Criteria Verification}: Every backlog item was checked against its predefined acceptance criteria. Only items meeting all criteria were marked as "Done."
        \item \textbf{Integration Testing}: All new endpoints and features were tested against the existing system to detect regressions. Postman collections were executed to verify API contracts.
        \item \textbf{Sprint Retrospective}: The team reflected on the sprint process — what went well, what could be improved — and identified actionable improvements for the next sprint.
    \end{spacing}
\end{itemize}

\subsection{Feedback Loop}
A continuous feedback mechanism was embedded throughout the development lifecycle:

\renewcommand{\labelitemi}{\tiny$\bullet$}
\begin{itemize}[leftmargin=2cm, topsep=4pt]
    \begin{spacing}{1.25}
        \item \textbf{Stakeholder demos}: At each sprint review, stakeholders interacted with the live system and provided direct feedback. This early and frequent exposure allowed the team to correct misunderstandings and adjust priorities before investing further effort.
        \item \textbf{Backlog refinement}: Feedback from each sprint was used to update and re-prioritize the Product Backlog. New requirements emerged during demos were added as new user stories, while lower-priority items were deferred.
        \item \textbf{Adaptive planning}: The Sprint Backlog for each iteration was adjusted based on lessons learned from the previous sprint. For example, the document upload feature was moved from Sprint 3 to Sprint 2 after stakeholders emphasized its importance during the Sprint 1 review.
        \item \textbf{Quality gates}: Automated tests and manual verification checkpoints prevented defects from accumulating. Any issue discovered during validation was logged as a bug in the backlog and scheduled for the next sprint.
    \end{spacing}
\end{itemize}

\begin{table}[H]
\centering
\begin{tabular}{|c|p{3.5cm}|p{4.5cm}|p{3.5cm}|}
\hline
\rowcolor{gray!20}
\textbf{Sprint} & \textbf{Main Objective} & \textbf{Delivered Features} & \textbf{Validation} \\
\hline
Sprint 1 & Authentication \& User Management & JWT login, RBAC, role management, base dashboards, change password & Stakeholder demo + Postman API tests \\
\hline
Sprint 2 & Reception Module & Candidate registration, year filtering, document upload, automated welcome email & Live receptionist workflow demo \\
\hline
Sprint 3 & Quiz Module & Quiz creation, candidate assignment, quiz interface, scoring, results & End-to-end quiz simulation with mock candidate \\
\hline
Sprint 4 & AI \& Assignment Module & Project ranking, candidate/supervisor matching, final assignment workflow & Supervisor assignment validation with sample data \\
\hline
\end{tabular}
\caption{SIPMS sprint planning with validation}
\end{table}

\section{CONCLUSION}
This chapter established the organizational context of the SIPMS project by presenting the host organization CLINISYS and analysing the critical limitations of its manual internship management system — including paper-based registration, unstructured document storage, subjective project assignment, and non-automated communication. In response to these shortcomings, we introduced the Smart Internship and Project Management System (SIPMS), a full-stack web platform built with React, Spring Boot, MySQL, and AI capabilities. The chapter also described the Scrum agile methodology adopted for development, organised into four sprints covering authentication, reception, quizzes, and AI-assisted assignment. These foundations provide the context for the detailed requirements analysis presented in the next chapter.
